\scanpage{41}
\margin{40.}
Soit $X$ un\deleted{\unsure{p|e}} preschema \uline{integre}, soit $R$ son corps des fonctions rationnelles (anneau local de l'element generique). Pour tout $x\quad X$, \deleted{$\unsure{O}$} l'anneau local $\underline O_x$ est integre et a $R$ comme corps des fractions, donc $\underline O_x$ s'identifie a un sous-anneau de $R$. Si $f\quad R$, on dit que $f$ est \uline{definie en} $x$ si $f\quad\underline O_x$, et l'ensemble des $x\quad X$ ou $f$ est definie est appele le \uline{domaine de definition} de $f$, note $D(f)$ (?). Il est evidemment \uline{ouvert}. Si $U\quad X$ est ouvert, on a $\mathrm I(U,\underline O_X)=\quad\underline O_x\quad(x\quad U)$

\uline{Proposition} 2 Soient $X$ un preschema integre de corps des fractions $R$. Pour que $X$ soit un schema, il faut et il suffit que pour tout couple de points $x,y\quad X$, $\unsure{\underline E_x|\underline O_x}$ et $\underline O_y$ apparentes implique $x=y$.

Supposons cette condition verifiee, montrons que $X$ est separe. Soient $U$ et $V$ deux ouverts affines de $X$, $A$ et $B$ leurs anneaux, identifies a des sous-anneaux de $R$. L'hypothese implique que $A$ et $B$ sont non apparentes (cf prop.1), \deleted{\unsure{donc $S(A)\quad S(B)$}} et $U$ ($V$) s'identifie a $S(A)$ ($S(B)$), alors $U\quad V$ s'identifie a $S(A)\quad S(B)=S(C)$, ou $C$ est l'\deleted{\unsure{anneau}} l'anneau engendre par $A$ et $B$ dans $R$. D'ailleurs, $C$ est identique a l'intersection des $Q\quad S(C)$, donc s'identifie a $\mathrm I(U\quad V,\underline O_X)$. Soit $W=U\quad V$, c'est un preschema tel que, posant $C=\mathrm I(W,\underline O_W)$, le morphisme naturel $W\quad\operatorname{Spec}(C)$ deduit de $C\quad\mathrm I(W,\underline O_W)$ soit bijectif (du point de vue ensembliste). Il s'ensuit facilement ici que c'est un homeomorphisme, car tout ferme de $W$ est la trace de son adherence dans $U$, donc defini par un ideal de $A$ et a fortiori par un ideal de \deleted{\unsure{xx}} $C$, donc est l'image reciproque d'un ferme de $\operatorname{Spec}(C)$. On voit de meme qu'on a en fait un \uline{isomorphisme} de preschemas, i.e. des isomorphismes pour les anneaux locaux. Donc $U\quad V$ est affine, d'autre part son anneau $C$ est engendre par les anneaux $A$ et $B$ de $U$ et $V$. Cela prouve que $X$ est separe (refernce).

Inversement, supposons $X$ separe, soient $x,y\quad X$ tels que $\underline O_x$ et $\underline O_y$ soient apparentes. Soit $U$ ($V$) un voisinage affine de $x$ ($y$) d'anneau $A$ ($B$ ), alors $U\quad V$ est affine d'anneau l'anneau $C$ engendre par $A$ et $B$ (meme reference). \deleted{En vertu de prop.1, dont \unsure{il}} Alors $x$\inserted{($y$)} correspond a un ideal premier $\underline p$ ($\underline q$) de $A$ ($B$), et comme \deleted{$\unsure{A_p}$} $A_{\underline p}$ et $B_{\underline q}$ sont apparentes il existe un ideal premier $\underline r$ de $C$ induisant $\underline p$ et $\underline q$. Alors $\underline r$ correspond a un point $z$ de $U\quad V$, et on a evidemment $x=z$ et $y=z$, d'ou $x=y$.

\uline{Corollaire} 1 Si $X$ est un schema integre, l'application $x\quad\underline O_x$ est un monomorphisme, i.e. si $x$ et $y$ sont des points distincts de $X$, alors il existe un fonction rationnelle definie en l'un et non en l'autre.

De façon plus precise:

\uline{Corollaire} 2 Soit $X$ un schema integre, soient $x,y$ deux points de $X$. Pour qu'on ait $x\quad\bar y$ i.e. $\bar x\quad\bar y$, il faut et il suffit que $\underline O_x\quad\underline O_y$, ie.
